% SPDX-License-Identifier: Apache-2.0
\section{\code{txhdl\_parts}}
\label{sec:r-parts}

What the language provides as hardware, written in its own lowerable
subset and checked as every unit is. A unit's channel port is one
side of a channel: the lowering makes of it a valid, a ready and the
data, and the testbench that checks the unit plays the channel's
other side from the trace. The channel itself, the elastic buffer of
two between two units, is not in either unit, so two lowered units
wired port to port would have no buffer between them, and their
handshakes would form a combinational loop, which a placer refuses.
\code{Buffer<W>} is the channel as hardware, one per channel of a run
wherever two lowered units meet, on a board or in a larger netlist.

Its rules are the runtime's, from Section~\ref{sec:r-comp}, with no
freedom taken. A take moves the tail into the head. A push goes into
the head if it is empty after that, else into the tail. \code{ready}
on the sending side is the tail being empty as the edge left it,
\code{valid} on the receiving side the head being full, and the data
is the head; all three are registers, so two units on a channel have
no combinational path between them. That the part is the runtime's
channel bit for bit is checked by the examples
document~\cite{examples}, which drives one with a pattern of offers
and takes, does the same on a runtime channel, and asserts after
every cycle that the two agree; the part's netlist is then checked
against that trace under both simulators.

\lstinputlisting[caption={\code{lib/parts/src/buffer.rs}.},label={lst:r-buffer}]{lib/parts/src/buffer.rs}
